\subsection{The case of Hilbert spaces}

In this section, we consider Hilbert spaces over K = R or C. We denote by \(\langle x_{1}|x_{2}\rangle\) the scalar product of two vectors \(x_{1}\) and \(x_{2}\) in a Hilbert space E.

If E is a complex Hilbert space, the Banach algebra \(\mathcal{L}(\mathrm{E})\) equipped with the involution \(u\mapsto u^{*}\) is a stellar algebra (example 1 of I, p.~102). In particular, if \(u\in \mathcal{L}(\mathrm{E})\) , one has \(\varrho (u^{*}) = \varrho (u)\) and \(\mathrm{Sp}(u^{*}) = \overline{\mathrm{Sp}(u)}\) .

Let \(u \in \mathcal{L}(\mathrm{E})\) be a normal endomorphism and let \(\lambda \in C\) . The eigenspace of u corresponding to \(\lambda\) coincides with the eigenspace of the adjoint \(u^{*}\) corresponding to \(\overline{\lambda}\) (EVT, V, p.~43, cor. of prop. 8). However, these spaces do not coincide in general if u is not normal (exercise 3 of I, p.~188).

Let \(\lambda\) and \(\mu\) be complex numbers such that \(\lambda \neq \mu\) . Let x be an eigenvector of u corresponding to \(\lambda\) and let y be an eigenvector of u corresponding to \(\mu\) . Then, \(u^{*}(x) = \overline{\lambda}x\) , hence

\[
\mu \langle x | y \rangle = \langle x | u (y) \rangle = \langle u ^ {*} (x) | y \rangle = \langle \overline {{\lambda}} x | y \rangle = \lambda \langle x | y \rangle .
\]

Consequently, \(\langle x|y\rangle = 0\) : the eigenspaces of u are pairwise orthogonal. Moreover, for every \(\lambda \in C\) , the eigenspace of u relative to \(\lambda\) coincides with the primary subspace of u relative to \(\lambda\) , that is (LIE, VII, §1, no. 1) the union over \(k \in N\) of the kernels of \((u - \lambda1_{\mathrm{E}})^{k}\) (EVT, V, p.~43, cor. of prop. 8).

Lemma 2. --- Let E be a complex Hilbert space and let u be a normal endomorphism of E. Let \((\mathrm{E}_{i})_{i\in\mathrm{I}}\) be a finite family of closed subspaces of E, stable under u and pairwise orthogonal, such that \(E=\bigoplus_{i\in I}E_{i}\) . For every \(i\in I\) , let \(u_{i}\) be the endomorphism of \(E_{i}\) induced by u. We have \(\mathrm{Sp}(u)=\bigcup_{i\in I}\mathrm{Sp}(u_{i})\) , and for every \(f\in\mathcal{C}(\mathrm{Sp}(u))\) , the endomorphism \(f(u)\) stabilizes the spaces \(E_{i}\) , and \(f(u)\) coincides with \(f(u_{i})\) on \(E_{i}\) .

The proof follows that of lemma 1 of I, p.~128, using remark 6 of I, p.~110 and prop. 8 of I, p.~112.

PROPOSITION 4. --- Let E be a complex Hilbert space and let u be a normal endomorphism of E. For every function \(f \in \mathcal{C}(\mathrm{Sp}(u))\) and every \(\lambda \in \mathbf{C}\) , one has

\[
\operatorname{Ker} (u - \lambda 1 _ {\mathrm{E}}) \subset \operatorname{Ker} (f (u) - f (\lambda) 1 _ {\mathrm{E}}).
\]

The proof is analogous to that of Proposition 1 of I, p.~128; let us take up its arguments again. The algebra A introduced in loc. cit. is here a unital stellar subalgebra of \(\mathcal{L}(\mathrm{E})\) (EVT, V, p.~43, cor.) . It therefore contains the unital stellar subalgebra B generated by \(u\) , which is commutative. The map \(\chi: \mathrm{B} \to \mathbf{C}\) which associates with \(v\) the eigenvalue of \(v\) relative to \(x\) is a character of B such that \(\chi(u) = \lambda\) . For every \(f \in \mathcal{C}(\mathrm{Sp}(u))\) , one has \(f(u) \in \mathrm{B}\) and \(\chi(f(u)) = f(\chi(u)) = f(\lambda)\) by prop. 8 of I, p.~112, whence the assertion.

Lemma 3. --- Let E be a Hilbert space and let \(p \in \mathcal{L}(\mathrm{E})\) be a projector. The following assertions are equivalent:

\begin{enumerate}
\def\labelenumi{(\roman{enumi})}
\item
  The projector p is an orthogonal projector, that is, \(\operatorname{Ker}(p)=\operatorname{Im}(p)^{\circ}\) (EVT, V, p.~13);
\item
  The projector p is Hermitian;
\item
  The projector p is normal;
\item
  We have \(\operatorname{Ker}(p) \subset \operatorname{Ker}(p^*)\) ;
\item
  We have \(\operatorname{Im}(p) \subset \operatorname{Im}(p^*)\) ;
\item
  The projector p is positive;
\item
  We have \(\| p\| \leqslant 1\) .
\end{enumerate}

Let us first recall that \(\operatorname{Ker}(p^{*}) = \operatorname{Im}(p)^{\circ}\) and \(\operatorname{Ker}(p) = \operatorname{Im}(p^{*})^{\circ}\) (EVT, V, p.~41, prop. 4). Moreover, the image of p (resp. \(p^{*}\) ) is closed, since it coincides with the kernel of 1 - p (resp. \(1 - p^{*}\) ). Hence we have

\[
\operatorname{Im} (p) = \operatorname{Ker} (p ^ {*}) ^ {\circ}, \qquad \operatorname{Im} (p ^ {*}) = \operatorname{Ker} (p) ^ {\circ}.\tag{1}
\]

\begin{enumerate}
\def\labelenumi{(\roman{enumi})}
\item
  \(\Longrightarrow\) (ii): \(p^{*}\) is a projector with kernel \(\operatorname{Im}(p)^{\circ} = \operatorname{Ker}(p)\) and whose image is \(\operatorname{Im}(p^{*})^{\circ} = \operatorname{Ker}(p) = \operatorname{Im}(p)^{\circ}\) ; therefore \(p^{*} = p\) .
\item
  \(\Longrightarrow\) (iii): since every Hermitian endomorphism is normal.
\item
  \(\Longrightarrow\) (iv): since \(p\) is normal, we have \(\| p(x)\|^2 = \| p^*(x)\|^2\) for all \(x\) in E (EVT, V, p.~43, prop. 7), whence the required inclusion.
\item
  \(\Longrightarrow\) (v) follows from equalities (1) above.
\item
  \(\Longrightarrow\) (vi): for all \(x \in E\) , one has \(p(x) \in \operatorname{Im}(p^{*})\) , and consequently \(\langle p(x)|x\rangle = \langle p^{*}(p(x))|x\rangle = \|p(x)\|^{2} \geqslant 0\) .
\item
  \(\Longrightarrow\) (vii): let \(x \in E\) and \(y = x - p(x) \in \operatorname{Ker}(p)\) . For every \(t \in R\) , by hypothesis
\end{enumerate}

\[
\langle x + t y | p (x) \rangle = \langle x + t y | p (x + t y) \rangle \geqslant 0,
\]

which is possible only if \(\langle y|p(x)\rangle = 0\) . But then

\[
\| p (x) \| ^ {2} = \langle x | p (x) \rangle \leqslant \| x \| \| p (x) \|,
\]

and therefore \(\|p\|\leqslant1\) .

\begin{enumerate}
\def\labelenumi{(\roman{enumi})}
\setcounter{enumi}{6}
\tightlist
\item
  \(\Longrightarrow\) (i): let \(y \in \operatorname{Im}(p)\) ; let \(z\) be the orthogonal projection of \(y\) on \(\operatorname{Ker}(p)^{\circ}\) and set \(x = y - z \in \operatorname{Ker}(p)\) . We have \(p(z) = p(y) = y\) , hence \(\|y\| \leqslant \|z\|\) by hypothesis. But, since \(x\) and \(z\) are orthogonal, we have \(\|y\|^2 = \|x\|^2 + \|z\|^2\) , whence \(\|x\| = 0\) , that is, \(y = z\) . Thus, \(\operatorname{Im}(p) \subset \operatorname{Ker}(p)^{\circ}\) . Since moreover \(\|p^*\| = \|p\| \leqslant 1\) , we likewise have \(\operatorname{Im}(p^*) \subset \operatorname{Ker}(p^*)^{\circ}\) which provides the reverse inclusion by (1).
\end{enumerate}

PROPOSITION 5. --- Let E be a complex Hilbert space and let u be a normal endomorphism of E.

\begin{enumerate}
\def\labelenumi{\alph{enumi})}
\item
  For every open and closed subset H of the spectrum of u, the spectral projector \(e_{\mathrm{H}}(u)\) is an orthogonal projector whose kernel is the image of the spectral projector \(e_{\mathrm{Sp}(u)-\mathrm{H}}(u)\) ;
\item
  If \(H_{1}\) and \(H_{2}\) are disjoint, open, and closed subsets of the spectrum of u, then the spectral subspaces \(E_{H_{1}}\) and \(E_{H_{2}}\) are orthogonal;
\item
  If \(\lambda \in \mathbf{C}\) is an isolated point of the spectrum of \(u\) , then \(\lambda\) is an eigenvalue of \(u\) and the image of the spectral projector \(e_{\lambda}(u)\) is the eigenspace of \(u\) corresponding to \(\lambda\) .
\end{enumerate}

We prove \(a\) ). Since the holomorphic functional calculus is compatible with the continuous functional calculus (I, p.~111, cor. 1), we have \(e_{\mathrm{H}}(u) = \varphi_{\mathrm{H}}(u)\) , where \(\varphi_{\mathrm{H}} \in \mathcal{C}(\mathrm{Sp}(u))\) is the characteristic function of H. This implies \(e_{\mathrm{H}}(u)^{*} = \overline{\varphi}_{\mathrm{H}}(u) = \varphi_{\mathrm{H}}(u) = e_{\mathrm{H}}(u)\) , hence \(e_{\mathrm{H}}(u)\) is an orthogonal projector (Lemma 3, (ii)). Its kernel is the image of the projector \(1 - e_{\mathrm{H}}(u) = e_{\mathrm{Sp}(u) - \mathrm{H}}(u)\) .

We prove \(b\) ). The characteristic functions \(\varphi_{\mathrm{H}_{1}}\) and \(\varphi_{\mathrm{H}_{2}}\) of \(\mathrm{H}_{1}\) and \(\mathrm{H}_{2}\) on \(\mathrm{Sp}(u)\) are continuous and their product is zero, which implies \(e_{\mathrm{H}_{1}}(u)\circ e_{\mathrm{H}_{2}}(u)=e_{\mathrm{H}_{2}}(u)\circ e_{\mathrm{H}_{1}}(u)=0\) . The inclusions \(\mathrm{E}_{\mathrm{H}_{2}}(u)\subset\mathrm{E}_{\mathrm{H}_{1}}(u)^{\circ}\) and \(\mathrm{E}_{\mathrm{H}_{1}}(u)\subset\mathrm{E}_{\mathrm{H}_{2}}(u)^{\circ}\) result from them.

Let us finally prove the assertion \(c\) ). The characteristic function \(\varphi_{\lambda}\) of \(\{\lambda\}\) is continuous and nonzero on \(\mathrm{Sp}(u)\) ; it satisfies \((z - \lambda)\varphi_{\lambda}(z) = 0\) for all \(z \in \mathrm{Sp}(u)\) . We therefore have \((u - \lambda 1_{\mathrm{E}})\varphi_{\lambda}(u) = 0\) . The nonzero image of \(\varphi_{\lambda}(u)\) is therefore contained in the eigenspace of \(u\) corresponding to \(\lambda\) . Since we have \(e_{\lambda}(u) = \varphi_{\lambda}(u)\) and since the image of \(e_{\lambda}(u)\) contains the eigenspace of \(u\) corresponding to \(\lambda\) , the assertion follows from this.

Lemma 4. — Let E be a Hilbert space and let u be a normal endomorphism of E. Let F be a closed subspace of E containing a total set of eigenvectors of u. Then \(\mathrm{F}^{\circ}\) is stable under u, and the endomorphism \(\widetilde{u}\) of \(\mathrm{F}^{\circ}\) induced by u is normal.

Since \(u\) is normal, every eigenvector of \(u\) is also an eigenvector of \(u^{*}\) (EVT, V, p.~43, cor.). The hypothesis therefore implies that F is stable under \(u\) and \(u^{*}\) . By EVT, V, p.~41, prop. 4 (ii), we therefore have \(u(\mathrm{F}^{\circ}) \subset \mathrm{F}^{\circ}\) and \(u^{*}(\mathrm{F}^{\circ}) \subset \mathrm{F}^{\circ}\) . It follows that the adjoint of \(\widetilde{u}\) is the endomorphism of \(\mathrm{F}^{\circ}\) induced by \(u^{*}\) . Since \(u\) is normal, the endomorphism \(\widetilde{u}\) is normal.

